\chapter{Implementation}
\label{ch:implementation}

\section{The from-scratch constraint}

Both models are implemented from first principles. The precise meaning of that constraint is
important, and is stated here before the code is presented.

PyTorch is used for three things: multidimensional array arithmetic, reverse-mode automatic
differentiation, and the optimiser. The modules \code{nn.Linear}, \code{nn.Embedding} and
\code{nn.Dropout} are used, but only as containers for learnable parameters --- an
\code{nn.Linear} is an affine map, not a sequence-modelling primitive.

Every operation specific to sequence modelling is written out. Table~\ref{tab:scratch} lists the
correspondence between the theory of Chapter~\ref{ch:theory} and the implementation.

\begin{table}[ht]
\centering
\caption{Components implemented from first principles, with the equations they realise.}
\label{tab:scratch}
\begin{tabular}{lll}
\toprule
\textbf{Component} & \textbf{Implemented as} & \textbf{Equations} \\
\midrule
LSTM gating               & \code{LSTMCellScratch}                & \ref{eq:lstm_i}--\ref{eq:lstm_h} \\
Recurrence, bidirectional & \code{LSTMLayerScratch}               & Sec.~\ref{sec:vanishing} \\
Recurrent classifier      & \code{BiLSTMClassifier}               & Sec.~\ref{sec:pooling} \\
Positional encoding       & \code{sinusoidal\us{}positional\us{}encoding} & Ch.~\ref{ch:theory} \\
Layer normalisation       & \code{LayerNormScratch}               & Ch.~\ref{ch:theory} \\
Scaled dot-product attention & \code{MultiHeadSelfAttentionScratch} & \ref{eq:attention} \\
Multi-head split and merge   & \code{MultiHeadSelfAttentionScratch} & \ref{eq:mha} \\
Encoder block             & \code{EncoderBlockScratch}            & Ch.~\ref{ch:theory} \\
Transformer classifier    & \code{TransformerClassifier}          & --- \\
\bottomrule
\end{tabular}
\end{table}

The following library primitives are \textbf{not used anywhere in either model}: \code{nn.LSTM},
\code{nn.GRU}, \code{nn.RNN}, \code{nn.MultiheadAttention}, \code{nn.TransformerEncoder},
\code{nn.TransformerEncoderLayer}, \code{nn.LayerNorm} and
\code{F.scaled\us{}dot\us{}product\us{}attention}. Two of these --- \code{nn.LSTM} and \code{F.layer\us{}norm}
--- appear exactly once each, inside the verification suite of
Chapter~\ref{ch:verification}, where they serve as numerical references.

\section{Model A: the bidirectional LSTM}

\subsection{The cell}

\code{LSTMCellScratch} implements Equations~\ref{eq:lstm_i}--\ref{eq:lstm_h} directly. The four
gate pre-activations are computed as a single fused product and split with \code{chunk(4)}.

{\small
\begin{verbatim}
class LSTMCellScratch(nn.Module):
    def __init__(self, d_in, d_hidden):
        super().__init__()
        self.d_hidden = d_hidden
        self.W = nn.Linear(d_in, 4 * d_hidden, bias=True)
        self.U = nn.Linear(d_hidden, 4 * d_hidden, bias=False)
        self.reset_parameters()

    def reset_parameters(self):
        k = 1.0 / math.sqrt(self.d_hidden)
        for p in self.parameters():
            nn.init.uniform_(p, -k, k)

    def forward(self, x_t, h_prev, c_prev):
        gates = self.W(x_t) + self.U(h_prev)        # (B, 4H) in one matmul
        i, f, g, o = gates.chunk(4, dim=-1)         # gate order: i, f, g, o
        i, f, o = torch.sigmoid(i), torch.sigmoid(f), torch.sigmoid(o)
        g = torch.tanh(g)
        c_t = f * c_prev + i * g
        h_t = o * torch.tanh(c_t)
        return h_t, c_t
\end{verbatim}}

Two design decisions deserve comment.

\keypoint{Gate ordering.} The chunks are ordered $i, f, g, o$. This ordering is arbitrary in
principle --- any consistent assignment trains equally well --- but it is chosen to match
PyTorch's internal convention exactly. That correspondence is what permits verification check~1 of
Chapter~\ref{ch:verification} to copy the reference implementation's weight matrices directly into
this cell and demand numerical equality. Choosing a different ordering would have made the
strongest available correctness test impossible.

\keypoint{Bias placement.} The reference implementation carries two bias vectors, one associated
with each of the input and hidden projections; their sum is mathematically equivalent to a single
bias. This implementation uses one bias on \code{W} and none on \code{U}. The verification check
accounts for this by copying the \emph{sum} of the two reference biases.

\subsection{The recurrence, and the masking problem}

\code{LSTMLayerScratch} unrolls the cell over the time dimension. Because batched sequences are
padded to a common length, this is the point at which a from-scratch implementation is most
likely to go subtly wrong.

{\small
\begin{verbatim}
class LSTMLayerScratch(nn.Module):
    def __init__(self, d_in, d_hidden, reverse=False):
        super().__init__()
        self.cell = LSTMCellScratch(d_in, d_hidden)
        self.d_hidden, self.reverse = d_hidden, reverse

    def forward(self, x, mask):
        B, T, _ = x.shape
        h = x.new_zeros(B, self.d_hidden)
        c = x.new_zeros(B, self.d_hidden)
        steps = range(T - 1, -1, -1) if self.reverse else range(T)
        outputs = []
        for t in steps:
            m = mask[:, t].unsqueeze(1)                   # (B, 1)
            h_new, c_new = self.cell(x[:, t], h, c)
            h = m * h_new + (1.0 - m) * h                 # padded step = no-op
            c = m * c_new + (1.0 - m) * c
            outputs.append(h)
        if self.reverse:
            outputs.reverse()                             # back to position order
        return torch.stack(outputs, dim=1), h
\end{verbatim}}

\keypoint{Problem 1: state corruption by padding.} The naive implementation applies the cell at
every timestep unconditionally. Padded positions then continue to drive the recurrence: the hidden
state at the nominal end of the sequence reflects an arbitrary number of pad-token updates rather
than the true end of the document. Because padding length varies across a batch, the magnitude of
this corruption varies per example.

The fix is to gate the state update by the mask,
\begin{equation}
h_t = m_t \odot \tilde{h}_t + (1 - m_t) \odot h_{t-1},
\end{equation}
and identically for $c_t$, where $m_t \in \{0,1\}$ indicates a real token. A padded step is then an
exact no-op: the state is carried forward unchanged.

\keypoint{Problem 2: the backward direction's starting point.} The backward recurrence must begin
at each sequence's \emph{own} final token, not at index 63. Sequences in a batch have different
true lengths, so there is no single starting index.

The implementation resolves this without any per-example index arithmetic, by exploiting the
masking just described. Padding sits at the right of each sequence, and masked steps are exact
no-ops. Iterating from $t = T-1$ down to $t = 0$ therefore passes through the padded region
performing no updates, and the recurrence effectively begins at each sequence's last real token.
The returned final state is the correct per-sequence endpoint in both directions.

This equivalence is an argument, not an observation, and arguments of this kind are exactly what
should be checked empirically. Verification checks~6 and~7 do so.

\subsection{The classifier}

{\small
\begin{verbatim}
class BiLSTMClassifier(nn.Module):
    def forward(self, ids, lengths=None):
        mask = (ids != PAD_ID).float()
        x = self.embedding(ids)
        out_f, h_f = self.forward_layer(x, mask)
        out_b, h_b = self.backward_layer(x, mask)
        seq = torch.cat([out_f, out_b], dim=-1)               # (B, T, 2H)
        denom = mask.sum(dim=1, keepdim=True).clamp(min=1.0)
        mean_pooled = (seq * mask.unsqueeze(-1)).sum(dim=1) / denom
        features = torch.cat([mean_pooled, h_f, h_b], dim=-1) # (B, 4H)
        return self.classifier(self.dropout(features))
\end{verbatim}}

The pooling follows Section~\ref{sec:pooling}: the masked mean over all positions is concatenated
with the final state of each direction, giving $4 d_h$ features. Note that the mean is computed
with an explicitly masked numerator and a denominator equal to the true length --- dividing by the
padded length would systematically shrink the representation of short documents.

\section{Model B: the Transformer encoder}

\subsection{Positional encoding and layer normalisation}

{\small
\begin{verbatim}
def sinusoidal_positional_encoding(max_len, d_model):
    position = torch.arange(max_len, dtype=torch.float32).unsqueeze(1)
    even_i   = torch.arange(0, d_model, 2, dtype=torch.float32)
    div_term = torch.exp(-math.log(10000.0) * even_i / d_model)
    pe = torch.zeros(max_len, d_model)
    pe[:, 0::2] = torch.sin(position * div_term)
    pe[:, 1::2] = torch.cos(position * div_term)
    return pe

class LayerNormScratch(nn.Module):
    def forward(self, x):
        mu  = x.mean(dim=-1, keepdim=True)
        var = x.var(dim=-1, keepdim=True, unbiased=False)
        return self.gamma * (x - mu) / torch.sqrt(var + self.eps) + self.beta
\end{verbatim}}

The \code{unbiased=False} argument is not incidental. The reference implementation normalises by
$d$ rather than $d-1$, and using the unbiased estimator here would produce a small but systematic
discrepancy. Verification check~2 confirms agreement to $5\times10^{-7}$.

\subsection{Multi-head self-attention}

{\small
\begin{verbatim}
class MultiHeadSelfAttentionScratch(nn.Module):
    def _split_heads(self, t):
        B, T, _ = t.shape
        return t.view(B, T, self.n_heads, self.d_k).transpose(1, 2)   # (B,h,T,d_k)

    def forward(self, x, key_mask=None, return_attention=False):
        B, T, _ = x.shape
        q = self._split_heads(self.W_q(x))
        k = self._split_heads(self.W_k(x))
        v = self._split_heads(self.W_v(x))

        scores = (q @ k.transpose(-2, -1)) / math.sqrt(self.d_k)      # (B,h,T,T)
        if key_mask is not None:
            scores = scores.masked_fill(key_mask[:, None, None, :] == 0, -1e9)

        attention = torch.softmax(scores, dim=-1)
        context = self.dropout(attention) @ v
        context = context.transpose(1, 2).reshape(B, T, self.n_heads * self.d_k)
        return self.W_o(context)
\end{verbatim}}

Three points are worth drawing out.

\keypoint{The head reshape.} The projections are computed once at full width and then reshaped
into $h$ heads. The sequence \code{view} then \code{transpose} is compact but opaque: it is easy to
write a variant that runs without error while silently interleaving the head dimensions, which
would mix features across heads. Verification check~3 targets this specifically by comparing
against a reference that slices head columns explicitly.

\keypoint{The mask broadcast.} The mask has shape $(B, T)$ and is indexed to $(B, 1, 1, T)$,
broadcasting across heads and across query positions. The semantics are that every query, in every
head, ignores padded \emph{keys}.

\keypoint{The finite mask constant.} The mask is applied as $-10^{9}$ rather than $-\infty$, for
the numerical reason given in Section~\ref{sec:masking}: a fully-masked row would otherwise
produce \code{NaN} and destroy the gradients of the entire batch.

\subsection{Encoder block and classifier}

{\small
\begin{verbatim}
class EncoderBlockScratch(nn.Module):
    def forward(self, x, key_mask=None):
        x = x + self.dropout(self.attention(self.norm1(x), key_mask))
        h = self.norm2(x)
        h = self.ff_out(self.dropout(F.gelu(self.ff_in(h))))
        return x + self.dropout(h)
\end{verbatim}}

The pre-norm ordering is visible directly: normalisation is applied to the \emph{input} of each
sublayer, and the residual stream itself is never normalised, leaving an unobstructed gradient
path from output to input.

The classifier embeds, scales by $\sqrt{d_{\text{model}}}$, adds the positional encoding, applies
the encoder blocks, applies a final normalisation, and reduces to a single vector by masked mean
pooling before the linear head.

\section{Training harness}

Both models are trained by a single function, so that the training budget is identical by
construction rather than by careful duplication.

\begin{table}[ht]
\centering
\caption{Training configuration. All values except the learning rate are shared.}
\label{tab:training_config}
\begin{tabular}{lll}
\toprule
\textbf{Setting} & \textbf{BiLSTM} & \textbf{Transformer} \\
\midrule
Optimiser              & \multicolumn{2}{l}{AdamW \cite{kingma2015,loshchilov2019}} \\
Weight decay           & \multicolumn{2}{l}{0.01} \\
Batch size             & \multicolumn{2}{l}{128} \\
Maximum epochs         & \multicolumn{2}{l}{8} \\
Gradient clipping      & \multicolumn{2}{l}{1.0 (global norm)} \\
Early stopping         & \multicolumn{2}{l}{validation macro-$F_1$, patience 2} \\
Checkpoint selection   & \multicolumn{2}{l}{best validation macro-$F_1$} \\
Random seed            & \multicolumn{2}{l}{42} \\
\midrule
Learning rate          & $10^{-3}$ & $3\times10^{-4}$ \\
\bottomrule
\end{tabular}
\end{table}

Gradient clipping at unit global norm addresses the exploding-gradient half of the pathology
analysed in Section~\ref{sec:vanishing}. Early stopping on macro-$F_1$ rather than
accuracy ensures that a model degrading on a single class is not selected.

\section{Architecture configuration and parameter budget}

\begin{table}[ht]
\centering
\caption{Architectural hyperparameters.}
\label{tab:arch_config}
\begin{tabular}{llll}
\toprule
\textbf{BiLSTM} & & \textbf{Transformer} & \\
\midrule
Embedding dimension & 128 & Model dimension $d_{\text{model}}$ & 128 \\
Hidden dimension    & 128 & Attention heads $h$               & 4 \\
Directions          & 2   & Encoder layers                    & 2 \\
Pooled features     & 512 & Feed-forward dimension $d_{ff}$   & 256 \\
Dropout             & 0.30 & Dropout                          & 0.10 \\
\bottomrule
\end{tabular}
\end{table}

These sizes were chosen to equalise non-embedding capacity, which is the quantity that must be
controlled for the comparison to be about architecture.

\begin{table}[ht]
\centering
\caption[Parameter budget]{Parameter budget. The embedding table is identical in size for both
models and dominates each total; the non-embedding column is therefore the meaningful measure of
architectural capacity.}
\label{tab:params}
\begin{tabular}{lrrr}
\toprule
\textbf{Model} & \textbf{Total} & \textbf{Embedding} & \textbf{Non-embedding} \\
\midrule
BiLSTM      & 4{,}612{,}612 & 4{,}347{,}392 & 265{,}220 \\
Transformer & 4{,}613{,}124 & 4{,}347{,}392 & 265{,}732 \\
\midrule
Difference  & 512 & 0 & 512 \ (0.19\%) \\
\bottomrule
\end{tabular}
\end{table}

The embedding table accounts for approximately 94\% of each model's parameters. Reporting only
total parameter counts would therefore obscure the comparison almost entirely --- the two models
would appear to differ by 0.01\%, which is true but uninformative. The non-embedding figures differ
by 512 parameters out of roughly 265{,}000, which is close enough that no observed performance
difference can reasonably be attributed to capacity.
